\documentclass[11pt]{article}
\usepackage[T1]{fontenc}
\usepackage[utf8]{inputenc}
\usepackage{lmodern,amsmath,amssymb,amsthm}
\usepackage[margin=1in]{geometry}
\usepackage{microtype}
\usepackage[colorlinks=true,linkcolor=blue,citecolor=blue,urlcolor=blue]{hyperref}
\hypersetup{pdftitle={Local Hessian integrability for planar Monge--Ampere inequalities},pdfauthor={Kaiwen Jin},pdfsubject={Planar Alexandrov solutions and local Hessian integrability}}
\newtheorem{theorem}{Theorem}[section]
\newtheorem{proposition}{Proposition}[section]
\newtheorem{lemma}{Lemma}[section]
\theoremstyle{remark}
\newtheorem{remark}{Remark}[section]
\numberwithin{equation}{section}
\newcommand{\R}{\mathbb R}
\newcommand{\C}{\mathbb C}
\newcommand{\D}{\mathbb D}
\newcommand{\Lip}{\operatorname{Lip}}
\newcommand{\SO}{\operatorname{SO}}
\newcommand{\tr}{\operatorname{tr}}
\newcommand{\supp}{\operatorname{supp}}
\title{Local Hessian integrability for planar Monge--Amp\`ere inequalities}
\author{Kaiwen Jin\\[2pt]\small\texttt{yui.ui.siki@gmail.com}}
\date{October 15, 2026}
\begin{document}
\maketitle

\begin{abstract}
Let \(0<\lambda<1\) and let \(u\) be a convex Alexandrov solution on a bounded
convex planar domain, with Monge--Amp\`ere density between \(\lambda\) and
\(\lambda^{-1}\). We prove that its full distributional Hessian belongs locally
to strong \(L^{p_\lambda}\), where
\(p_\lambda=(1+\lambda^2)/(1-\lambda^2)\). This implies the corresponding weak
endpoint estimate posed in the Monge--Amp\`ere regularity literature.
The proof uses a change of variables with a Lipschitz inverse and a local
weighted scalar Lipschitz projection. An energy estimate on the set where
the projection differs from the original map controls the cutoff error
and yields the required negative moment.
\end{abstract}

\section{Introduction}\label{sec:intro}

The Monge--Amp\`ere equation controls the determinant of a convex potential's
Hessian. Even when that determinant is bounded above and below, its two
eigenvalues can be very different. Estimates on the whole Hessian therefore
contain information about anisotropy that the determinant alone does not
display. This distinction is important in the regularity of convex potentials
and their gradient maps; the connection with optimal transportation is
described in the survey of De Philippis and Figalli~\cite{DPF}.

For strictly convex Alexandrov solutions with two-sided density bounds,
De Philippis, Figalli and Savin~\cite[Theorem~1.1]{DFS} proved local
\(W^{2,1+\varepsilon}\) regularity for some positive \(\varepsilon\) depending
on the density bounds and dimension. In the plane, strict convexity follows
from the positive lower bound. The local Hessian is thus a matrix field, but
the quantitative dependence of its integrability exponent remains a separate
question. De Philippis and Figalli~\cite[\S5.2(2)]{DPF} formulate the weak
endpoint exponent
\[
p_\lambda=\frac{1+\lambda^2}{1-\lambda^2},
\qquad 0<\lambda<1,
\]
under the bounds \(\lambda\le\det D^2u\le\lambda^{-1}\).
We establish the following strong local estimate at this exponent.

\begin{theorem}\label{thm:main}
Let \(0<\lambda<1\), let \(\Omega\subset\R^2\) be a bounded open convex domain,
and let \(u:\Omega\to\R\) be finite and convex. Suppose that its Alexandrov
Monge--Amp\`ere measure satisfies
\begin{equation}\label{eq:MA}
\lambda|E|\le\mu_u(E)=|\partial u(E)|\le\lambda^{-1}|E|
\quad\text{for every Borel }E\subset\Omega.
\end{equation}
Then the full distributional Hessian is a measurable matrix field \(H\), and
for every \(\Omega'\Subset\Omega\),
\begin{equation}\label{eq:mainstrong}
\int_{\Omega'}(1+|H|)^{p_\lambda}\,dx<\infty,
\qquad p_\lambda=\frac{1+\lambda^2}{1-\lambda^2}.
\end{equation}
In particular,
\begin{equation}\label{eq:mainweak}
\sup_{t>0}t^{p_\lambda}
|\{x\in\Omega':|H(x)|>t\}|<\infty.
\end{equation}
\end{theorem}

Here \(\partial u(x)\) is the convex subgradient, \(\partial u(E)\) is the union
of those subgradients over \(E\), and \(|E|\) denotes planar Lebesgue measure.
The matrix norm is the operator norm. The conclusion concerns the entire
distributional Hessian, including the assertion that no singular Hessian
measure remains. Constants may depend on the fixed solution, \(\lambda\),
and the compact interior set. Chebyshev's inequality gives
\eqref{eq:mainweak} from \eqref{eq:mainstrong}.

We use the separating function and scalar projection ideas of Agazzi
et al.~\cite[Lemmas~3--4]{Agazzi}. The local argument here is weighted.
The periodic quasiconvexity theorem in that paper is not invoked on a local
domain. The cone and reflection-map calculations follow the geometry
developed by De Philippis, Guerra and Tione~\cite[Lemmas~6.3 and~6.9]{DPGT};
we give the transformations directly for the original Alexandrov class.

The change of variables \(y=x+Du(x)\) turns the Hessian into the derivative
of a 1-Lipschitz map. Large Hessian eigenvalues correspond to a small distance
of that derivative from a circle of rotations. The inverse of each rotated
difference is quasiconformal, which gives a quantitative distance deficit
for the Lipschitz map on fixed interior disks. A scalar Lipschitz projection then has a uniformly small
error. Its constraint is saturated wherever that error is nonzero, and the
resulting energy estimate controls the distance on precisely that set.
A sufficiently high power of a smooth cutoff gives an integrable error when
the separating inequality is integrated over its scale.

Section~\ref{sec:coordinates} establishes the exact measure transformation
and the derivative cone. Section~\ref{sec:projection} proves the local weighted
estimate. Section~\ref{sec:moments} integrates it and returns to the Hessian.
Throughout, vector norms are Euclidean and matrix norms are operator norms.
We identify \(\R^2\) with \(\C\), and \(C\) denotes a finite local constant
independent of the small projection scale and the integration cutoff.

\section{Coordinates for Alexandrov solutions}\label{sec:coordinates}

We first justify the change of variables without assuming a smooth solution.
The Radon--Nikodym theorem~\cite[Theorem~2]{TaoRN} and
\eqref{eq:MA} give
\begin{equation}\label{eq:density}
\mu_u=f\,dx,\qquad \lambda\le f\le\lambda^{-1}\quad\text{a.e.}
\end{equation}
Planar strict convexity and local \(C^{1,\alpha}\) regularity follow from
the results recorded in~\cite[p.~18 and Theorem~2.14]{DPF}.
The normalized \(W^{2,1+\varepsilon}\) theorem~\cite[Theorem~1.1]{DFS}
applies locally as follows. At \(x_0\), subtract the supporting affine
function \(\ell\) and choose a small compact section
\(S=\{u-\ell<h\}\Subset\Omega\).
The potential \(w=u-\ell-h\) is zero on its boundary and has minimum \(-h\)
at \(x_0\). John normalization
\cite[Lemma~2.15]{DPF} gives an affine map \(A(x)=Lx+c\), \(\det L>0\),
with \(B_1\subset A(S)\subset B_2\). In dimension two,
\((\det L)w\circ A^{-1}\) has the same density bounds: its subgradient area
and its domain area both acquire the factor \(\det L\).
The minimum point belongs to the half-height set in the theorem of
\cite{DFS}. Pulling back the resulting weak derivatives gives
\[
u\in C^1_{\rm loc}(\Omega)\cap W^{2,1}_{\rm loc}(\Omega).
\]
Write \(H=D^2u\) for this full weak matrix field. It is symmetric and
positive semidefinite almost everywhere by convexity.

\begin{proposition}\label{prop:measure}
Set
\[
\psi(x)=u(x)+\tfrac12|x|^2,\qquad
Y(x)=x+Du(x),\qquad U=Y(\Omega).
\]
Then \(Y:\Omega\to U\) is a homeomorphism onto an open set, with
1-Lipschitz inverse \(F\). Both maps send null sets to null sets. Moreover,
\begin{equation}\label{eq:transport}
\det H=f,\qquad
\rho=1+\tr H+\det H,\qquad
|Y(E)|=\int_E\rho\,dx
\end{equation}
for every Borel \(E\subset\Omega\). Almost everywhere in \(U\),
\begin{equation}\label{eq:inverse}
DF(y)=(I+H(F(y)))^{-1}.
\end{equation}
\end{proposition}

\begin{proof}
Convex mollification on compact subdomains and continuity of
Monge--Amp\`ere measures~\cite[Proposition~2.4]{DPF} give, for fixed \(t\ge0\),
\begin{equation}\label{eq:mixed}
\mu_{u+t|x|^2/2}=\mu_u+t\Delta u+t^2\,dx.
\end{equation}
Indeed the smooth identity is
\(\det(D^2u_j+tI)=\det D^2u_j+t\Delta u_j+t^2\).
The two convex measures converge weakly on compact tests, and the
Laplacians converge distributionally. This proves \eqref{eq:mixed}.
Consequently, with
\(\rho_0=1+\tr H+f\), we have \(\mu_\psi=\rho_0\,dx\).

Monotonicity of \(Du\) gives
\[
(Y(x)-Y(z))\cdot(x-z)\ge |x-z|^2.
\]
Thus \(Y\) is injective and \(F\) is 1-Lipschitz. To see that \(U\) is open,
take \(\overline B_a(x_0)\subset\Omega\).
On its sphere, strong convexity gives
\[
\psi(x)-y\cdot x-\psi(x_0)+y\cdot x_0
\ge a^2/2-|y-Y(x_0)|a.
\]
For \(|y-Y(x_0)|<a/2\), the minimum on this ball is interior,
so \(D\psi=y\) somewhere in the ball. This also proves continuity of the
inverse on the image. Since \(\psi\in C^1\), its subgradient is \(Y\), and
\begin{equation}\label{eq:transport0}
|Y(E)|=\int_E\rho_0\,dx.
\end{equation}
Hence \(Y\) sends null sets to null sets. The Lipschitz map \(F\) does too,
by the area formula~\cite[\S2, Theorem~1]{Jerrard}.

The convex Legendre supremum
\(\psi^*(y)=\sup_{x\in\Omega}(x\cdot y-\psi(x))\) is finite:
\(\Omega\) is bounded and \(\psi\) has an affine lower bound.
For \(y\in U\), its unique maximizer is \(F(y)\). The inequalities
\[
F(y)\cdot h\le \psi^*(y+h)-\psi^*(y)\le F(y+h)\cdot h
\]
show that \(D\psi^*=F\) on \(U\), since the remainder is at most \(|h|^2\).
Thus \(G=DF\), defined almost everywhere by Rademacher's theorem
\cite[Theorem~2.2.4]{Tao}, is symmetric and positive semidefinite.
Let \(B\) consist of its singular-derivative set and the null set where
the derivative does not exist. The area formula gives \(|F(B)|=0\);
\eqref{eq:transport0} then gives \(|B|=|Y(F(B))|=0\).
Therefore \(G\) is nonsingular almost everywhere.

At such a point \(y\), the expansion
\(F(y+h)-F(y)=Gh+o(|h|)\) implies a lower bound \(c|h|\).
Continuity of the inverse converts this expansion into differentiability
of \(Y\) at \(F(y)\), with derivative \(G^{-1}\).
The exceptional set has null image under \(F\). The ordinary derivative
of this continuous \(W^{1,1}_{\rm loc}\) map agrees almost everywhere
with its weak derivative \(I+H\).
One can verify this on almost every coordinate line: Fubini gives the
one-dimensional weak derivative, its integral is an absolutely continuous
representative, continuity identifies that representative, and Lebesgue
differentiation identifies its ordinary derivative. These integration
and differentiation statements are given in~\cite[Theorems~1.6.11--12
and~1.7.21]{Tao}. The null-image property of \(Y\) now gives
\eqref{eq:inverse} in image coordinates.

Applying the area formula to \(F\) on \(Y(E)\), and using the weighted form
of \eqref{eq:transport0}, yields
\[
|E|=\int_{Y(E)}\det G\,dy
=\int_E\frac{\rho_0}{\det(I+H)}\,dx.
\]
Equality on every Borel set implies
\(\rho_0=\det(I+H)=1+\tr H+\det H\), so \(\det H=f\).
This proves \eqref{eq:transport}.
To apply the globally stated area formula, each component of \(F\) may be
extended to \(\R^2\) by the infimum Lipschitz extension. Its general
multiplicity formula is then restricted to sets in \(U\), where \(F\)
is injective. The extension need not be injective outside \(U\).
\end{proof}

Define
\begin{equation}\label{eq:Wv}
W(y)=y-2F(y),\qquad
v(y)=A_0W(y),\qquad A_0=\operatorname{diag}(1,-1).
\end{equation}
For \(x,z\in\Omega\), monotonicity gives
\[
|W(Y(x))-W(Y(z))|^2
=|Y(x)-Y(z)|^2-4(Du(x)-Du(z))\cdot(x-z)
\le |Y(x)-Y(z)|^2.
\]
Thus \(W\) and \(v\) are 1-Lipschitz. Their derivatives satisfy
\begin{equation}\label{eq:cayley}
DW(Y(x))=(H-I)(H+I)^{-1}.
\end{equation}

For a real matrix \(M\), write its complex action as
\(Mz=a_+(M)z+a_-(M)\bar z\). Its operator norm and determinant are
\[
|M|=|a_+(M)|+|a_-(M)|,\qquad
\det M=|a_+(M)|^2-|a_-(M)|^2.
\]
Since \(DW\) is symmetric, the derivative at almost every \(y\in U\) acts by
\begin{equation}\label{eq:complex}
Dv(y)\xi=b(y)\xi+a(y)\bar\xi\quad(\xi\in\C),\qquad
b=b_1-ib_2,\quad a\in\R.
\end{equation}
We suppress the base point in the coefficient functions below. Put
\begin{equation}\label{eq:Kd}
K=\lambda^{-1},\qquad k=\frac{K-1}{K+1},\qquad d=1-|b|.
\end{equation}

\begin{proposition}\label{prop:cone}
Almost everywhere in \(U\),
\begin{equation}\label{eq:cone}
0<d\le1,\qquad |a|\le kd,\qquad
|Dv|\le1-(1-k)d.
\end{equation}
If \(h_1\ge h_2>0\) are the eigenvalues of \(H\), then
\begin{equation}\label{eq:compare}
\frac{2}{(1+k)(1+h_1)}\le d(Y(x))
\le\frac{2K}{1+h_1},
\qquad 1+h_1\le\rho\le C_K(1+h_1).
\end{equation}
For every rotation \(Q\in\SO(2)\), the map
\(g_Q(y)=v(y)-Qy\) is a \(K\)-quasiconformal homeomorphism onto its image.
\end{proposition}

\begin{proof}
The eigenvalues of \(DW\) are \((h_i-1)/(h_i+1)\in(-1,1)\);
therefore \(|b|+|a|<1\). The inverse Cayley formula gives
\[
\det H=\frac{(1+a)^2-|b|^2}{(1-a)^2-|b|^2}.
\]
The determinant band is equivalent to
\(2|a|\le k(1+a^2-|b|^2)\). Hence
\[
|a|\le\frac{k(1-|b|)(1+|b|)}{2-k|a|}
\le k(1-|b|),
\]
where the last inequality follows from \(|b|+k|a|\le1\).
This proves \eqref{eq:cone}.
The matrix computation is the operator-norm cone calculation in
\cite[Lemma~6.3]{DPGT}.

The operator distance of \(DW\) from symmetric orthogonal reflections is
\[
d+|a|=\max\left\{\frac{2}{1+h_1},\frac{2h_2}{1+h_2}\right\}.
\]
To check the formula, a reflection has eigenvalues \(1,-1\).
Its two eigendirections give lower bounds \(1-\alpha_1\) and
\(1+\alpha_2\), where \(\alpha_1\ge\alpha_2\) are the eigenvalues of \(DW\);
aligning these directions with those of \(DW\) attains their maximum.
Equivalently, left multiplication by \(A_0\) turns this distance into
\(\min_{|Q|=1}(|b-Q|+|a|)=d+|a|\).
Since \(h_2=f/h_1\),
\[
\frac{h_2(1+h_1)}{1+h_2}
=\frac{f(1+h_1)}{h_1+f}\le\max\{f,1\}\le K.
\]
Together with \(d\le d+|a|\le(1+k)d\), this proves the distance
comparison in \eqref{eq:compare}. The density comparison follows from
\(\rho=1+h_1+h_2+f\), \(h_2\le h_1\), and \(f\le K\).

The complex coefficients of \(Dg_Q\) are \(b-Q,a\), and
\(|a|\le kd\le k|b-Q|\).
Thus \(|Dg_Q|^2\le K\det Dg_Q\), with positive determinant.
To establish homeomorphism, set \(S=A_0Q\), a symmetric reflection.
Choose orthonormal \(e,e^\perp\) with \(Se=e\), \(Se^\perp=-e^\perp\).
In original coordinates,
\[
W(Y(x))-SY(x)
=-2(x\cdot e)e+2(Du(x)\cdot e^\perp)e^\perp.
\]
Equal outputs fix \(x\cdot e\). On that line, \(Du\cdot e^\perp\)
is strictly increasing by differentiable strict convexity, so the inputs
coincide. The map is also open: in a small rectangle about any point,
the derivative values at the upper and lower ends of its central vertical
segment strictly surround the central value. Continuity preserves a
smaller such interval for all nearby horizontal coordinates, and
the intermediate value theorem fills an open output rectangle.
This proves homeomorphism after composition with \(F\) and \(A_0\).
It is the \(C^1\) fiber argument underlying~\cite[Lemma~6.9]{DPGT}.
The analytic quasiconformal convention
\cite[\S2.1]{DPGT} is precisely
\(|Dg_Q|^2\le K\det Dg_Q\) for a \(W^{1,2}_{\rm loc}\) homeomorphism,
so its \(K\) agrees with \eqref{eq:Kd}.
\end{proof}

\section{A local weighted projection estimate}\label{sec:projection}

\subsection{A distance deficit on interior disks}

Fix concentric disks
\(\overline B_r(y_0)\subset B_R(y_0)\Subset U\), with \(0<r<R\).
All constants in this section may depend on these disks and the fixed map.

\begin{lemma}\label{lem:short}
There is \(c>0\) such that, for all \(y,z\in\overline B_r(y_0)\),
\begin{equation}\label{eq:short}
|v(y)-v(z)|\le |y-z|-c|y-z|^K.
\end{equation}
\end{lemma}

\begin{proof}
Joint continuity, compactness of \(\SO(2)\), and the injectivity from
Proposition~\ref{prop:cone} give
\[
\ell=\min_{\substack{Q\in\SO(2),\,y\in\overline B_r\\z\in\partial B_R}}
|g_Q(y)-g_Q(z)|>0.
\]
The image \(V_Q=g_Q(B_R)\) is bounded and simply connected, and
\(\partial V_Q=g_Q(\partial B_R)\). These facts follow by pulling loops
back into the disk and using the homeomorphism on its compact closure.
Choose a Riemann map \(\chi_Q:V_Q\to\D\) with
\(\chi_Q(g_Q(y_0))=0\), using~\cite[p.~1]{SiuR}.
The disk self-homeomorphism
\(T_Q(z)=\chi_Q(g_Q(y_0+Rz))\) and its inverse are
\(K\)-quasiconformal and fix zero. Inversion and conformal composition
preserve \(K\); see~\cite[Theorem~4.16 and the following paragraph]{Calegari}.
Mori's theorem~\cite[Theorem~1.1]{Mori}, in its planar normalized form,
states
\[
|T_Q^{-1}(\xi)-T_Q^{-1}(\xi')|
\le16|\xi-\xi'|^{1/K}.
\]

The set \(g_Q(\overline B_r)\) is at boundary distance at least \(\ell\)
in \(V_Q\). If two such image points have distance \(\sigma<\ell/2\),
their connecting segment is inside \(V_Q\), at distance at least
\(\ell/2\) from its boundary. Cauchy's derivative formula
\cite[p.~5]{SiuC}, on circles of radius \(\ell/4\), gives
\(|\chi_Q'|\le4/\ell\) along this segment. Mori therefore gives
\[
|y-z|\le16R(4\sigma/\ell)^{1/K}.
\]
If \(\sigma\ge\ell/2\), use \(|y-z|\le2r\).
Thus, uniformly in \(Q\),
\begin{equation}\label{eq:inverseholder}
|y-z|\le C|g_Q(y)-g_Q(z)|^{1/K}.
\end{equation}
For a pair with \(v(y)\ne v(z)\), choose \(Q\) taking the direction of
\(y-z\) to that of \(v(y)-v(z)\).
Because \(v\) is 1-Lipschitz, its image difference under \(g_Q\) then has
length \(|y-z|-|v(y)-v(z)|\). Raising \eqref{eq:inverseholder} to
power \(K\) proves \eqref{eq:short}.
When the vector difference vanishes, decrease \(c\) to at most
\((2r)^{1-K}\); the same inequality holds.
\end{proof}

\subsection{The scalar projection}

Let \(f_1=v_1\). Choose \(\zeta\in C_c^\infty(B_r)\), \(0\le\zeta\le1\),
equal to one on \(B_{r/2}\). Fix an integer \(m>K+1\) and set
\(\eta=\zeta^m\). The high power will be used to control \(D\eta\).
For \(0<\epsilon<1/2\), put \(s=1-\epsilon\).
Choose a scalar \(s\)-Lipschitz minimizer \(w\) of
\[
\int_{B_r}\eta(f_1-w)^2,
\]
and set \(h=f_1-w\). The construction adapts the projection in
\cite[Lemma~4]{Agazzi} to this local weight.

\begin{lemma}\label{lem:projection}
The minimizer can be chosen so that
\begin{equation}\label{eq:residual}
\|h\|_\infty\le C\epsilon^{K/(K-1)}.
\end{equation}
On \(E=\{h\ne0\}\cap\{\eta>0\}\), \(|Dw|=s\) almost everywhere.
On \(\{h=0\}\), \(Dh=0\) almost everywhere. In addition,
\begin{equation}\label{eq:translation}
\int\eta Dw\cdot Dh\ge-\int hD\eta\cdot Dw.
\end{equation}
\end{lemma}

\begin{proof}
Existence follows from compactness of bounded equi-Lipschitz scalar
functions on the closed disk. The positive mass of \(\eta\) and a bounded
objective bound one value of a minimizing sequence, and the Lipschitz
constraint bounds all its values. A uniformly convergent subsequence has
an admissible limit with the limiting objective.
The variational inequality is
\begin{equation}\label{eq:VI}
\int\eta h(\xi-w)\le0
\quad\text{for every scalar \(s\)-Lipschitz competitor \(\xi\).}
\end{equation}

Define the two \(s\)-Lipschitz envelopes
\[
f_s^-(y)=\inf_{z\in\overline B_r}(f_1(z)+s|y-z|),
\qquad
f_s^+(y)=\sup_{z\in\overline B_r}(f_1(z)-s|y-z|).
\]
They satisfy \(f_s^-\le f_1\le f_s^+\).
By Lemma~\ref{lem:short}, either deviation from \(f_1\) is bounded by
\(\sup_{t\ge0}(\epsilon t-ct^K)\le C\epsilon^{K/(K-1)}\).
Clipping a minimizer between these envelopes preserves its Lipschitz
constraint and does not increase its squared distance from \(f_1\).
This proves \eqref{eq:residual}.

At a differentiability point where \(h>0\), \(\eta>0\), and \(|Dw|<s\),
a small cone \(w(y_1)+t-s|y-y_1|\) lies below \(w\) on a sufficiently
small surrounding sphere and above it at the center.
Replace \(w\) by its maximum with that cone inside the sphere.
The functions agree at the boundary; splitting line segments there
shows that the modification remains \(s\)-Lipschitz on the original disk.
It increases \(w\) on an open subset where \(\eta h>0\), contradicting
\eqref{eq:VI}. The minimum with an upward cone handles \(h<0\).
This proves saturation almost everywhere. The zero-set assertion follows
at its density and differentiability points: a nonzero derivative would
produce a conical sector of nonzero values, contradicting density one.

Extend \(w\) to \(\R^2\) with the same Lipschitz bound.
Both translations \(w(\cdot+te_i)\) and \(w(\cdot-te_i)\), restricted
to the disk, are competitors in \eqref{eq:VI}.
With \(A=\eta h\), extended by zero, their summed inequality is
\[
\begin{aligned}
0&\ge\int A(y)\bigl(w(y+te_i)+w(y-te_i)-2w(y)\bigr)\,dy\\
&=-\int\bigl(A(y+te_i)-A(y)\bigr)
          \bigl(w(y+te_i)-w(y)\bigr)\,dy.
\end{aligned}
\]
Divide by \(t^2\). The first differences converge almost everywhere to
the respective first derivatives, and are bounded on a fixed compact set.
Dominated convergence gives \(\int\partial_iA\,\partial_iw\ge0\).
Summing over \(i\) proves \eqref{eq:translation}.
Only first differences were used; compact support controls the translations
without a boundary term.
\end{proof}

\subsection{The separating inequality and energy estimate}

For a real matrix \(M\), define
\begin{equation}\label{eq:Ls}
L_s(M)=
\begin{cases}
\det M,& |M|\le s,\\
2s|a_+(M)|-s^2,& |M|>s,
\end{cases}
\qquad D_\epsilon(M)=L_{1-\epsilon}(M)-\det M.
\end{equation}
Here \(a_+(M)\) is the coefficient of \(z\) in
\(Mz=a_+(M)z+a_-(M)\bar z\).
This is the rescaled separating function of~\cite[Lemma~3]{Agazzi}.
We record the pointwise calculation that will be used.

\begin{lemma}\label{lem:calibration}
If \(M_1,M_2\) are the rows of \(M\), then
\[
L_s(M)\ge \det(z,M_2)+z\cdot(M_1-z)
\]
whenever \(|z|=s\), or whenever \(z=M_1\) and \(|M_1|\le s\).
\end{lemma}

\begin{proof}
Let \(J=\left(\begin{smallmatrix}0&-1\\1&0\end{smallmatrix}\right)\).
For \(|z|=s\), the supremum of the right side is
\(s|M_1-JM_2|-s^2=2s|a_+(M)|-s^2\).
Write \(\alpha=|a_+(M)|\), \(\tau=|a_-(M)|\).
When \(|M|\le s\), the difference between \(\det M\) and this supremum is
\((s-\alpha)^2-\tau^2\ge0\), because \(s\ge\alpha+\tau\).
When \(|M|>s\), the definition already equals the supremum.
For \(z=M_1\), the right side is \(\det M\).
If \(|M|>s\) but \(|M_1|\le s\), then
\(|\alpha-\tau|\le|M_1|\le s<\alpha+\tau\);
hence \((s-\alpha)^2-\tau^2\le0\), so \(L_s(M)\ge\det M\)
also in this case.
\end{proof}

\begin{proposition}\label{prop:defect}
For the fixed disks and cutoff above, let
\[
\gamma=\frac1{K-1}-\frac1m,\qquad
\beta=\frac2{K^2-1}.
\]
Then \(\gamma>\beta\), and for all \(0<\epsilon<1/2\),
\begin{equation}\label{eq:defect}
\int\eta D_\epsilon(Dv)\,dy\ge-C\epsilon^{2+\gamma}.
\end{equation}
\end{proposition}

\begin{proof}
By Lemma~\ref{lem:projection}, \(Dw\) has length \(s\) on \(E\),
and equals \(Df_1\) where \(h=0\). Lemma~\ref{lem:calibration} therefore gives
\[
D_\epsilon(Dv)\ge-\det(Dh,Dv_2)+Dw\cdot Dh
\quad\text{on }\{\eta>0\}.
\]
The real-antilinear form \eqref{eq:complex} gives the rows
\[
Df_1=(b_1+a,b_2),\qquad Dv_2=(-b_2,b_1-a).
\]
Thus \(B=-JDv_2=Df_1-2ae_1\) is distributionally divergence free.
Testing against \(\eta h\) and using \eqref{eq:translation},
\begin{equation}
\begin{aligned}
\int\eta D_\epsilon(Dv)
&\ge \int hD\eta\cdot(Df_1-Dw)-2\int ha\,\partial_1\eta\\
&=-\frac12\int h^2\Delta\eta-2\int ha\,\partial_1\eta.
\end{aligned}
\label{eq:cancel}
\end{equation}
These integrations by parts use compactly supported Lipschitz tests and
bounded first derivatives. Smooth approximation of the tests justifies
them in the weak sense.

To estimate the last term, use saturation on its actual nonzero set.
The identity \(Df_1=Dw+Dh\) gives
\[
2\int\eta Dw\cdot Dh+\int\eta|Dh|^2
=\int_E\eta(|Df_1|^2-s^2).
\]
From \eqref{eq:cone}, \(|Df_1|\le1-(1-k)d\), while \(s^2\ge1-2\epsilon\).
Since \((1-t)^2\le1-t\) for \(0\le t\le1\),
\[
|Df_1|^2-s^2\le2\epsilon-(1-k)d.
\]
Combining this with \eqref{eq:translation} and \eqref{eq:residual} yields
\begin{equation}\label{eq:energy}
(1-k)\int_E\eta d+\int\eta|Dh|^2
\le2\epsilon\int_E\eta+2\|h\|_\infty\int|D\eta|
\le C\epsilon.
\end{equation}

The chosen power of the cutoff satisfies
\(|D\eta|\le C\eta^{1-1/m}\).
H\"older's inequality and \(d\le1\) now give
\[
\int_Ed|D\eta|
\le C\left(\int_E\eta d\right)^{1-1/m}
       \left(\int_Ed\right)^{1/m}
\le C\epsilon^{1-1/m}.
\]
Since \(|a|\le kd\) and \(\|h\|_\infty\le C\epsilon^{1+1/(K-1)}\),
the two terms in \eqref{eq:cancel} satisfy
\[
\left|\int ha\,\partial_1\eta\right|
\le C\epsilon^{2+1/(K-1)-1/m},
\qquad
\left|\int h^2\Delta\eta\right|
\le C\epsilon^{2+2/(K-1)}.
\]
This proves \eqref{eq:defect}. Finally,
\[
\gamma-\beta=\frac1{K+1}-\frac1m>0,
\]
by the choice \(m>K+1\).
\end{proof}

\section{Negative moments and Hessian integrability}\label{sec:moments}

We integrate \eqref{eq:defect} over its scale. The sign of that integral
must be controlled uniformly for every positive lower cutoff.

\begin{lemma}\label{lem:scalar}
Let \(0<d\le1\), \(0\le A\le kd\), and define
\[
P_\epsilon=
\begin{cases}
0,&\epsilon\le d-A,\\
A^2-(d-\epsilon)^2,&\epsilon>d-A.
\end{cases}
\]
For every \(0<q<(1-k)/k\), there are constants \(c_q>0,C_q<\infty\),
depending only on \(q,k\), such that
\begin{equation}\label{eq:scalar}
\int_{h_0}^{1/2}P_\epsilon\,\epsilon^{-q-3}\,d\epsilon
\le C_q-c_q\max\{d,h_0\}^{-q},
\qquad 0<h_0<1/2.
\end{equation}
\end{lemma}

\begin{proof}
Put \(b_0=A/d\in[0,k]\). Direct integration of the three monomials gives
\[
\begin{aligned}
I_0&=\int_{d-A}^{\infty}
[A^2-(d-\epsilon)^2]\epsilon^{-q-3}\,d\epsilon\\
&=\frac{2((q+1)b_0-1)}{q(q+1)(q+2)}
(1-b_0)^{-q-1}d^{-q}
\le-c_0d^{-q},
\end{aligned}
\]
where \(c_0=2(1-(q+1)k)/(q(q+1)(q+2))>0\).
All denominators are positive, since \(d-A\ge(1-k)d>0\).
If \(h_0\le d-A\), the continued integral from \(h_0\) equals \(I_0\).
If \(d-A<h_0\le d+A\), the removed part is nonnegative, so the remainder
is at most \(I_0\). In either case it is bounded above by
\(-c_0\max\{d,h_0\}^{-q}\).

If \(h_0>d+A\), the remaining integrand is negative.
For \(\epsilon\ge4h_0\), using \(A,d\le h_0\), we have
\[
A^2-(d-\epsilon)^2\le-\epsilon^2/2.
\]
Its tail is at most \(-(4h_0)^{-q}/(2q)\).
Finally, the absolute integral of the continued expression from \(1/2\)
to infinity is bounded uniformly in \(d,A\le1\), because
\(|A^2-(d-\epsilon)^2|\le C(1+\epsilon^2)\) and \(q>0\).
Removing that tail proves \eqref{eq:scalar}.
This continuation is only a scalar polynomial calculation.
\end{proof}

\begin{proof}[Proof of Theorem~\ref{thm:main}]
For the map \(v\) from Section~\ref{sec:coordinates}, the definition
\eqref{eq:Ls} gives, with \(A=|a|\),
\[
D_\epsilon(Dv)=0\quad(\epsilon\le d-A),\qquad
D_\epsilon(Dv)=A^2-(d-\epsilon)^2\quad(\epsilon>d-A).
\]
Take \(q=\beta=2/(K^2-1)\). Then
\[
0<q<\frac{1-k}{k}=\frac2{K-1},\qquad \gamma>q.
\]
By Proposition~\ref{prop:defect},
\[
\int_{h_0}^{1/2}
\left(\int\eta D_\epsilon(Dv)\,dy\right)\epsilon^{-q-3}\,d\epsilon
\ge-C\int_0^{1/2}\epsilon^{\gamma-q-1}\,d\epsilon>- \infty,
\]
uniformly in \(h_0>0\).
At each positive cutoff, the integrand is absolutely integrable on the
product of \(\supp\eta\) and \([h_0,1/2]\).
Fubini's theorem, in its absolute-integrability form
\cite[Theorem~1.7.21 and Corollary~1.7.23]{Tao}, and
Lemma~\ref{lem:scalar} therefore imply a uniform bound on
\(\int\eta\max\{d,h_0\}^{-\beta}\,dy\).
As \(h_0\downarrow0\), these nonnegative functions increase.
Monotone convergence~\cite[Theorem~1.4.44]{Tao} gives
\begin{equation}\label{eq:moment}
\int_{B_{r/2}}d^{-\beta}\,dy<\infty.
\end{equation}
No measurable choice of the projections as \(\epsilon\) varies is required:
the integrated expression is the explicit \(D_\epsilon(Dv)\).

For \(\Omega'\Subset\Omega\), the set \(Y(\overline{\Omega'})\) is compact
in \(U\). Finitely many inner disks of the kind above cover it, so
\(\int_{Y(\Omega')}d^{-\beta}\,dy<\infty\).
The weighted inverse area formula
\cite[\S2, Theorem~2]{Jerrard}, Proposition~\ref{prop:measure},
and \eqref{eq:compare} give
\[
\begin{aligned}
\int_{\Omega'}(1+|H|)^{1+\beta}\,dx
&=\int_{Y(\Omega')}
\frac{(1+h_1(F(y)))^{1+\beta}}{\rho(F(y))}\,dy\\
&\le C_K\int_{Y(\Omega')}d(y)^{-\beta}\,dy<\infty.
\end{aligned}
\]
Indeed \(\rho\ge1+h_1\) and \(d\le2K/(1+h_1)\).
Since \(1+\beta=(1+\lambda^2)/(1-\lambda^2)\), this is
\eqref{eq:mainstrong}. Finally,
\(t^{p_\lambda}\mathbf1_{\{|H|>t\}}\le|H|^{p_\lambda}\)
for every \(t>0\), proving \eqref{eq:mainweak}.
\end{proof}

\section*{AI-use statement}
OpenAI Codex was used with substantial involvement in developing the argument,
checking source applications, and preparing the exposition.
An AI reader and an AI referee, each in a fresh context, reviewed the same
initial anonymous manuscript. Their reports concern that snapshot; one author
revision followed with AI assistance.
Independent mathematical review is requested.

\begingroup
\small
\raggedright
\begin{thebibliography}{11}
\bibitem{Agazzi}
A.~Agazzi, K.~Astala, R.~Ba\~nuelos, G.~Bruno, G.~Cassese,
D.~Faraco, A.~Guerra, M.~Hormozi, B.~Kirchheim, A.~Koski,
V.~Kova\v{c}, J.~Kristensen, F.~Pasqualotto, I.~Prause,
E.~Saksman, V.~\v{S}ver\'ak, L.~Sz\'ekelyhidi, R.~Tione,
A.~Volberg and H.~Yang,
\emph{Sharp Higher Integrability Theory, Part I:
The \(L^p\)-norm of the Beurling--Ahlfors transform},
arXiv:2609.28311v2 (2026).
\url{https://arxiv.org/abs/2609.28311v2}.

\bibitem{Calegari}
D.~Calegari, \emph{Chapter 2: Hyperbolic Geometry},
author's manuscript, \S4.2, Definitions~4.13 and~4.15, Remark~4.14,
and Theorem~4.16.
\url{https://math.uchicago.edu/~dannyc/books/3manifolds/3_manifolds_chapter_2.pdf}.

\bibitem{DPF}
G.~De Philippis and A.~Figalli,
\emph{The Monge--Amp\`ere equation and its link to optimal transportation},
Bulletin of the American Mathematical Society \textbf{51} (2014), 527--580.
\url{https://doi.org/10.1090/S0273-0979-2014-01459-4}.
Page and theorem locators refer to the preprint arXiv:1310.6167v1,
\url{https://arxiv.org/abs/1310.6167v1}.

\bibitem{DFS}
G.~De Philippis, A.~Figalli and O.~Savin,
\emph{A note on interior \(W^{2,1+\varepsilon}\) estimates for the
Monge--Amp\`ere equation},
Mathematische Annalen \textbf{357} (2013), 11--22.
\url{https://doi.org/10.1007/s00208-012-0895-9};
preprint version \url{https://arxiv.org/abs/1202.5566v3}.

\bibitem{DPGT}
G.~De Philippis, A.~Guerra and R.~Tione,
\emph{Unique continuation for differential inclusions},
Annales de l'Institut Henri Poincar\'e C \textbf{43} (2026), 127--154.
\url{https://doi.org/10.4171/AIHPC/146}.

\bibitem{Mori}
R.~Fehlmann and M.~Vuorinen,
\emph{Mori's theorem for \(n\)-dimensional quasiconformal mappings},
Annales Academiae Scientiarum Fennicae, Series A I Mathematica
\textbf{13} (1988), 111--124.
\url{https://doi.org/10.5186/aasfm.1988.1304}.

\bibitem{Jerrard}
R.~L.~Jerrard, \emph{Area and coarea},
MAT1501 lecture notes, University of Toronto (2013), \S2, Theorems~1--2.
\url{https://www.math.toronto.edu/rjerrard/1501/notes/week5-6.pdf}.

\bibitem{SiuR}
Y.-T.~Siu, \emph{Riemann Mapping Theorem},
Math 113 lecture notes, Harvard University (2024).
\url{https://people.math.harvard.edu/~siu/math113/15_riemann_mapping_theorem.pdf}.

\bibitem{SiuC}
Y.-T.~Siu, \emph{Cauchy's Integral Formula and Power and Laurent Series Expansion},
Math 113 lecture notes, Harvard University (2024).
\url{https://people.math.harvard.edu/~siu/math113/5_cauchy_integral_formula_power_and_laurent_series_expansion.pdf}.

\bibitem{Tao}
T.~Tao, \emph{An Introduction to Measure Theory},
Graduate Studies in Mathematics \textbf{126},
American Mathematical Society (2011).
The cited theorem and corollary numbers refer to the linked preliminary
version made available by the author (2012):
\url{https://terrytao.wordpress.com/wp-content/uploads/2012/12/gsm-126-tao5-measure-book.pdf}.

\bibitem{TaoRN}
T.~Tao, \emph{245B, notes 1: Signed measures and the
Radon--Nikodym--Lebesgue theorem}, author's lecture notes (4 January 2009).
\url{https://terrytao.wordpress.com/2009/01/04/245b-notes-1-signed-measures-and-the-radon-nikodym-lebesgue-theorem/}.
\end{thebibliography}
\endgroup
\end{document}
